% B-diagnostics-impl.tex — Annex B: Diagnostic Catalogue
\chapter{Diagnostic Catalogue}
\label{annex:B-diagnostics-impl}
\index{diagnostics!catalogue}
\index{error codes}

\pnum
Every diagnostic emitted by the Lain compiler carries a bracketed code of
the form \texttt{[EXXX]}.  Each diagnostic is written to \texttt{stderr} and
causes a non-zero exit code.  This annex lists every code in numeric order
together with the triggering condition and the source location responsible.

\pnum
Diagnostics whose prefix is \texttt{[VRA]} are emitted by the bounds checker
(\srcref{src/sema/bounds.h}) and do not carry a structured \texttt{E-number};
they are documented separately in \cref{sec:bdiag-vra}.

% -------------------------------------------------------------------------
\section{Linearity and ownership errors (E001–E010)}
\label{sec:bdiag-linearity}
\index{linearity errors}
% -------------------------------------------------------------------------

\begin{longtable}{p{2.2cm}p{4cm}p{7.5cm}}
\caption{Linearity and ownership diagnostic codes.\label{tab:diag-linearity}} \\
\toprule
\textbf{Code} & \textbf{Summary} & \textbf{Triggering condition} \\
\midrule
\endfirsthead
\multicolumn{3}{l}{\small (continued)} \\
\toprule
\textbf{Code} & \textbf{Summary} & \textbf{Triggering condition} \\
\midrule
\endhead
\bottomrule
\endfoot

\diagcode{E001} &
Use after move &
An identifier whose \cname{LState} is \cname{LSTATE\_CONSUMED} is referenced.
Also raised when a borrow's root owner has been moved and the borrow is still
used (\srcref{src/sema/region.h}). \\

\diagcode{E002} &
Double move / field re-consume &
A \lain{mov} expression is applied to an already-consumed variable, or a
struct field whose \cname{FieldState} is already \cname{LSTATE\_CONSUMED} is
consumed a second time (\srcref{src/sema/linearity.h}). \\

\diagcode{E003} &
Unconsumed linear value &
A linear variable or field reaches a \lain{return} or the end of its scope
without having been consumed.  Also raised when a \cname{STMT\_EXPR} discards
a value of a \cname{mov} type (\srcref{src/sema/linearity.h}). \\

\diagcode{E004} &
Borrow aliasing violation &
A new borrow conflicts with an existing live borrow: (a) a shared borrow is
requested while a mutable borrow is active, (b) a mutable borrow is requested
while any borrow (shared or mutable) is active.  Also raised when an access is
attempted on a mutably-borrowed variable (\srcref{src/sema/region.h}). \\

\diagcode{E005} &
Use of uninitialized variable &
An identifier is accessed before any assignment has been recorded for it in
the \cname{LTable} (\srcref{src/sema/linearity.h}). \\

\diagcode{E006} &
Loop-level consume &
A linear variable defined outside a loop body is consumed inside the loop, or
a field of such a variable is consumed inside the loop.  This would cause
multiple consumption across iterations (\srcref{src/sema/linearity.h}). \\

\diagcode{E007} &
Missing \lain{mov} keyword &
A function parameter of \cname{MODE\_OWNED} type is passed without an explicit
\lain{mov} at the call site (\srcref{src/sema/linearity.h}). \\

\diagcode{E008} &
Move while borrowed &
A \lain{mov} expression is applied to a variable that still has at least one
active borrow, or the variable has partially-consumed fields and cannot be
moved as a whole (\srcref{src/sema/linearity.h}). \\

\diagcode{E009} &
Assignment to immutable &
An assignment target is a variable declared without \lain{var} (i.e.\
\cname{StmtVar.is\_mutable} is false) (\srcref{src/sema/resolve.h}). \\

\diagcode{E010} &
Dangling reference &
A mutable reference (\lain{var T}) is returned from a function and the
referenced variable is local to that function (lifetime analysis in both
\srcref{src/sema/resolve.h} and \srcref{src/sema/region.h}). \\

\end{longtable}

% -------------------------------------------------------------------------
\section{Purity and mutability errors (E011)}
\label{sec:bdiag-purity}
\index{purity errors}
% -------------------------------------------------------------------------

\begin{longtable}{p{2.2cm}p{4cm}p{7.5cm}}
\caption{Purity diagnostic codes.\label{tab:diag-purity}} \\
\toprule
\textbf{Code} & \textbf{Summary} & \textbf{Triggering condition} \\
\midrule
\endfirsthead
\multicolumn{3}{l}{\small (continued)} \\
\toprule
\textbf{Code} & \textbf{Summary} & \textbf{Triggering condition} \\
\midrule
\endhead
\bottomrule
\endfoot

\diagcode{E011} &
Purity violation &
A \lain{func} declaration violates one of three purity rules:
(1) it calls a \lain{proc} (side-effectful procedure);
(2) it assigns to a global variable;
(3) it participates in a mutual-recursion cycle with another \lain{func}
    (checked via the DFS in \srcref{src/sema.h});
(4) it contains an unbounded \lain{while} loop (one without a
    \lain{decreasing} measure).  All cases share the same code but produce
    distinct message text. \\

\end{longtable}

% -------------------------------------------------------------------------
\section{Type and call errors (E012–E015)}
\label{sec:bdiag-type}
\index{type errors}
% -------------------------------------------------------------------------

\begin{longtable}{p{2.2cm}p{4cm}p{7.5cm}}
\caption{Type and call diagnostic codes.\label{tab:diag-type}} \\
\toprule
\textbf{Code} & \textbf{Summary} & \textbf{Triggering condition} \\
\midrule
\endfirsthead
\multicolumn{3}{l}{\small (continued)} \\
\toprule
\textbf{Code} & \textbf{Summary} & \textbf{Triggering condition} \\
\midrule
\endhead
\bottomrule
\endfoot

\diagcode{E012} &
Type/call mismatch &
A catch-all code for multiple static type errors:
wrong argument count to a function;
struct initializer argument count mismatch;
struct field type mismatch in a positional call;
arithmetic binary operator applied to a struct or enum type;
generic function call missing a type argument;
unsafe pointer cast outside an \lain{unsafe} block;
type alias RHS does not evaluate to a type.
Each instance carries a distinct human-readable message. \\

\diagcode{E013} &
Name error &
Covers two sub-cases:
(1) variable redeclaration or shadowing — a new \cname{STMT\_VAR} uses a name
already bound in the current scope;
(2) assignment to an undeclared identifier;
(3) \lain{main} is declared as \lain{func} rather than \lain{proc}. \\

\diagcode{E014} &
Non-exhaustive match / comptime error &
Raised in two contexts:
(1) a \lain{case} expression or \lain{match} statement does not cover all
    variants of the scrutinee type (exhaustiveness check in
    \srcref{src/sema/typecheck.h});
(2) a \lain{comptime if} condition cannot be reduced to a compile-time
    constant;
(3) CTFE recursion depth exceeds \cname{COMPTIME\_MAX\_DEPTH = 256}
    (\srcref{src/sema/comptime.h}). \\

\diagcode{E015} &
Division by zero &
Inside a \lain{func}, VRA determines that the right-hand operand of a
\texttt{/} or \texttt{\%} expression may be zero (range includes 0, or range
is fully unknown) (\srcref{src/sema/typecheck.h}). \\

\end{longtable}

% -------------------------------------------------------------------------
\section{Termination-measure errors (E080–E082)}
\label{sec:bdiag-term}
\index{termination errors}
% -------------------------------------------------------------------------

\pnum
Codes E080–E082 are emitted by the bounded-while verifier embedded in
\srcref{src/sema.h}.  They apply only to \lain{while} loops that carry a
\lain{decreasing} measure clause.

\begin{longtable}{p{2.2cm}p{4cm}p{7.5cm}}
\caption{Termination diagnostic codes.\label{tab:diag-term}} \\
\toprule
\textbf{Code} & \textbf{Summary} & \textbf{Triggering condition} \\
\midrule
\endfirsthead
\multicolumn{3}{l}{\small (continued)} \\
\toprule
\textbf{Code} & \textbf{Summary} & \textbf{Triggering condition} \\
\midrule
\endhead
\bottomrule
\endfoot

\diagcode{E080} &
Measure non-negativity failure &
The verifier cannot prove that the termination measure is $\geq 0$ when the
loop condition is true.  Specifically, no non-negativity pattern matches: the
measure is not of the form $b - a$ with $b \geq a$ via an in-guard, not a
single variable with a proven $\geq 0$ range, and not reducible by and-chain
traversal. \\

\diagcode{E081} &
No variables in measure &
\cname{measure\_extract\_vars} returns an empty variable set.  This means the
measure expression contains no identifiers and cannot decrease as a function
of any program variable. \\

\diagcode{E082} &
Measure decrease not provable &
The verifier cannot confirm that the body strictly decreases the measure.
Either no assignment to any measure variable is found in the body, or all
found assignments have \cname{ASSIGN\_DIRECTION\_UNKNOWN} or
\cname{ASSIGN\_DIRECTION\_INCREASE} direction. \\

\end{longtable}

% -------------------------------------------------------------------------
\section{Parser errors (E100)}
\label{sec:bdiag-parser}
\index{parser errors}
% -------------------------------------------------------------------------

\begin{longtable}{p{2.2cm}p{4cm}p{7.5cm}}
\caption{Parser diagnostic codes.\label{tab:diag-parser}} \\
\toprule
\textbf{Code} & \textbf{Summary} & \textbf{Triggering condition} \\
\midrule
\endfirsthead
\multicolumn{3}{l}{\small (continued)} \\
\toprule
\textbf{Code} & \textbf{Summary} & \textbf{Triggering condition} \\
\midrule
\endhead
\bottomrule
\endfoot

\diagcode{E100} &
Parse error &
All parser errors share this code.  Specific messages include:
semicolons not allowed (Lain uses newlines as statement separators);
unexpected token at top level (an unrecognised declaration keyword);
unexpected token in expression (primary expression parse failed).
The parser calls \cname{parser\_error()} which prints the message and calls
\cname{exit(1)} (\srcref{src/parser/core.h}). \\

\end{longtable}

% -------------------------------------------------------------------------
\section{VRA bounds diagnostics}
\label{sec:bdiag-vra}
\index{VRA!diagnostics}
% -------------------------------------------------------------------------

\pnum
The Value-Range Analysis bounds checker (\srcref{src/sema/bounds.h}) does not
use the \texttt{[EXXX]} code scheme.  It prefixes messages with \texttt{[VRA]}
and emits one of four conditions:

\begin{longtable}{p{3.5cm}p{9.8cm}}
\caption{VRA bounds-check messages.\label{tab:diag-vra}} \\
\toprule
\textbf{Message text (prefix)} & \textbf{Condition} \\
\midrule
\endfirsthead
\multicolumn{2}{l}{\small (continued)} \\
\toprule
\textbf{Message text (prefix)} & \textbf{Condition} \\
\midrule
\endhead
\bottomrule
\endfoot

\texttt{index may be negative} &
The index expression's \cname{Range.min} is $< 0$. \\

\texttt{index N out of bounds} &
Both index and length are statically known and index $\geq$ length. \\

\texttt{index range unknown \ldots} &
The index range is unknown and the array length is known; a runtime
bounds check is emitted into the C output. \\

\texttt{cannot prove index < length} &
The index range is known but the array length is unknown; a runtime
bounds check is emitted. \\

\texttt{both index and length unknown} &
Neither bound is known at compile time; a runtime bounds check is
emitted.  This is the common case for dynamically-sized slices with
non-trivial index expressions. \\

\end{longtable}

\pnum
The first two conditions are hard errors (exit code 1).  The last three are
informational: they indicate that a runtime check has been inserted, not that
an error has been found.  However, the compiler may be configured to treat
unproved accesses as warnings in future versions; currently the distinction is
only visible in the diagnostic prefix.

% -------------------------------------------------------------------------
\section{Diagnostic format}
\label{sec:bdiag-format}
\index{diagnostics!format}
% -------------------------------------------------------------------------

\pnum
All structured diagnostics follow the format:

\begin{ccode}
[EXXX] Error Ln <line>, Col <col>: <message>
\end{ccode}

\pnum
For diagnostics that also invoke \cname{diagnostic\_show\_line()}, the output
continues with two additional lines:

\begin{ccode}
    <source line text>
    <caret string>
\end{ccode}

\noindent
where the caret string contains a \texttt{\^{}} aligned under the column
position.  The function \cname{diagnostic\_show\_line()} is declared in
\srcref{src/sema.h} and called at approximately 68 sites across the semantic
analysis passes.

\pnum
The VRA messages and the older borrow-checker messages in
\srcref{src/sema/region.h} do not yet include the \texttt{Ln}/\texttt{Col}
prefix on all call sites and do not invoke \cname{diagnostic\_show\_line()}.
Unifying these is a known area for future improvement.
